% Conclusion for the Nera draft (Oct 23, 2026). Numbers as in Sections 5-6 of this draft.
\section{Conclusion}
\label{sec:conclusion}

I build a life-cycle model of fertility and housing in which households
choose when to have children, how much space to occupy, whether to rent or
own, and how much to save. Children need space, the largest homes must be
bought with a down payment, and earnings are uninsurable. House prices and
rents are determined in equilibrium, and the children born today become
households a generation later. I calibrate the model to the United States in
2007, take the fall in fertility since then as given, and study housing
policies from the 2023 economy.

Mortgage credit for all households mainly moves births earlier. Young
households use it to buy the small homes they would otherwise rent, and the
completed fertility of the 2023 cohort falls by
1.1\%.\textcolor{red}{$^{\ddagger}$} Policies that lower the cost of a child
raise completed fertility: more housing supply, credit for parents, by 4.9\%,
and insurance against low earnings, by 14.0\%, against 7.8\% for cash. A
rebated property tax lowers the house price but raises rents, and a lower
price is not cheaper space. In the long run births per household return to
replacement, so the gains show up as more households and build up over a
generation. Policy raises births when it lowers the cost of a child rather
than the cost of a home, and housing policy for births must be judged by its
long-run effects.

Several features of the model matter for the size of these effects. The model
overstates homeownership among young households, which places too many of
them near the margin where credit acts. A common decline in the preference
for children overstates childlessness in 2023 and misses the postponement of
first births, and it makes the house price fall after 2007, while in the data
it rose. Children ever born by age 25 are below the data. Credit for parents
does not price default or mortgage-insurance costs, and the transfers are not
financed. Landlords and developers are absentee, so rental income and capital
gains on rental housing leave the economy.
